\documentclass[10pt]{article}
\usepackage[T1]{fontenc}
\usepackage[utf8]{inputenc}
\usepackage[french]{babel}
\usepackage{amsmath,amssymb,mathrsfs,mathpazo}
\usepackage{geometry}
\geometry{a4paper,margin=2cm}
\usepackage{graphicx,array,tabularx,tabularray,colortbl,multirow,multicol,booktabs,enumitem,pifont,ulem,fancybox,diagbox,cancel}
\usepackage{pstricks,pstricks-add,pst-plot,pst-tree,pst-node,pst-text,pst-func,pst-math,pst-bspline,pst-3dplot,pst-eucl,pst-solides3d}
\let\pscalculate\undefined
\usepackage{pst-bezier}
\usepackage{tikz,pgfplots,tkz-tab}
\usetikzlibrary{arrows,arrows.meta,patterns,calc,positioning,trees,decorations.markings,shapes.geometric,shapes.misc,shapes.arrows,intersections,decorations.pathmorphing,decorations.pathreplacing,backgrounds,shadows,babel}
\pgfplotsset{compat=1.18}
\usepgfplotslibrary{fillbetween}
\usepackage[np]{numprint}
\usepackage{listings,makecell,pas-tableur,twoopt}
\let\barre\undefined
\usepackage[pstricks]{tablvar}
\definecolor{darkgreen}{RGB}{0,100,0}
\providecommand{\red}{\color{red}}
\providecommand{\blue}{\color{blue}}
\providecommand{\green}{\color{green}}
\providecommand{\black}{\color{black}}
\providecommand{\white}{\color{white}}
\providecommand{\gray}{\color{gray}}
\providecommand{\cyan}{\color{cyan}}
\providecommand{\magenta}{\color{magenta}}
\providecommand{\yellow}{\color{yellow}}
\providecommand{\violet}{\color{violet}}
\providecommand{\orange}{\color{orange}}
\providecommand{\euro}{\texteuro}
\providecommand{\EUR}{\texteuro}
\providecommand{\ieme}{\textsuperscript{e}}
\providecommand{\degres}{\ensuremath{^\circ}}
\providecommand{\degree}{\ensuremath{^\circ}}
\providecommand{\vv}[1]{\overrightarrow{#1}}
\providecommand{\Coord}[1]{#1}
\providecommand{\decosix}{}
\providecommand{\starredbullet}{\ensuremath{\bullet}}
\pagestyle{empty}
\setlength{\parindent}{0pt}
\usepackage{adjustbox}\DeclareUnicodeCharacter{2236}{:}\DeclareUnicodeCharacter{2190}{\ensuremath{\leftarrow}}
\providecommand{\R}{\mathbb{R}}
\providecommand{\N}{\mathbb{N}}
\providecommand{\Z}{\mathbb{Z}}
\providecommand{\Q}{\mathbb{Q}}
\providecommand{\D}{\mathbb{D}}
\providecommand{\C}{\mathbb{C}}
\providecommand{\e}{{\mathrm{e}}}
\providecommand{\vect}[1]{\overrightarrow{#1}}
\providecommand{\vectt}[1]{\overrightarrow{\mathrm{#1}}}
\providecommand{\vvt}[1]{\overrightarrow{\mathrm{#1}}}
\providecommand{\np}[2][]{{#2~#1}}
\providecommand{\vv}[1]{{\overrightarrow{#1}}}
\renewcommand{\d}{{\mathrm{d}}}
\providecommand{\degres}{{^\circ}}
\providecommand{\degre}{{^\circ}}
\providecommand{\blue}{\color{blue}}
\providecommand{\red}{\color{red}}
\providecommand{\enskip}{\quad}
\providecommand{\llbracket}{\lbrack\!\lbrack}
\providecommand{\rrbracket}{\rbrack\!\rbrack}
\providecommand{\euro}{{\text{€}}}
\providecommand{\no}{{\text{n°}}}
\providecommand{\interval}[2]{{[#1;#2]}}
\providecommand{\intervFO}[2]{{[#1;#2[}}
\providecommand{\dotfill}{\ldots}

\providecommand{\ds}{\displaystyle}
\providecommand{\textding}[1]{#1}
\providecommand{\Oij}{\(\mathrm{O};\vec{\imath},\vec{\jmath}\)}
\providecommand{\Oijk}{\(\mathrm{O};\vec{\imath},\vec{\jmath},\vec{k}\)}
\providecommand{\Ouv}{\(\mathrm{O};\vec{u},\vec{v}\)}
\providecommand{\intervFF}[2]{\ensuremath{[#1;#2]}}
\renewcommand{\euro}{\ifmmode\text{\texteuro}\else\texteuro\fi}
\DeclareUnicodeCharacter{00B0}{\ifmmode^{\circ}\else\textdegree\fi}
\DeclareUnicodeCharacter{20AC}{\ifmmode\text{\texteuro}\else\texteuro\fi}
% Énoncés : Éducation nationale. Transcription et corrigés : APMEP, Denis Vergès et contributeurs. Découpage et classement : CoopMaths, Sébastien Lozano, Éric Elter et contributeurs. Adaptation de la présentation : Maths Explorer.
% CC BY-SA : https://coopmaths.fr/www/about/
% https://coopmaths.fr/alea/static/bac/2024/tex/bac_2024_11_sujet2_amsud_3.tex
\begin{document}
\noindent\textit{APMEP / CoopMaths. CC BY-SA. Présentation adaptée par Maths Explorer.}\par\medskip
\medskip

\textbf{Partie 1}

\medskip

On considère la fonction $f$ définie sur l'ensemble des nombres réels $\R$ par :
\[f(x) = \left(x^2 - 4\right)\e^{-x}.\]

On admet que la fonction $f$ est dérivable sur $\R$ et on note $f'$ sa fonction dérivée.

\medskip

\begin{enumerate}
\item Déterminer les limites de la fonction $f$ en $-\infty$ et en $+\infty$.
\item Justifier que pour tout réel $x,\, f'(x) = \left(-x^2 + 2x + 4\right)\e^{-x}$.
\item En déduire les variations de la fonction $f$ sur $\R$.
\end{enumerate}

\bigskip

\textbf{Partie 2}

\medskip

On considère la suite $(I_n)$ définie pour tout entier naturel $n$ par $I_n = \displaystyle\int_{-2}^0 x^n\e^{-x}\, \text{d}x$.

\medskip

\begin{enumerate}
\item Justifier que $I_0 = \e^2 - 1$.
\item En utilisant une intégration par partie, démontrer l'égalité :

\[I_{n+1} = (- 2)^{n+1} \text{e}^2 + (n + 1)I_n.\]

\item En déduire les valeurs exactes de $I_1$ et de $I_2$.
\end{enumerate}

\medskip

\textbf{Partie 3}

\medskip

\begin{minipage}{0.55\linewidth}
\begin{enumerate}
\item Déterminer le signe sur $\R$ de la fonction $f$ définie dans la partie 1.
\item On a représenté ci-contre la courbe $\mathcal{C}_f$ de la fonction $f$ dans un repère orthonormé \Oij.

Le domaine $D$ du plan hachuré ci-contre est délimité par la courbe $\mathcal{C}_f$, l'axe des abscisses et l'axe des ordonnées.

Calculer la valeur exacte, en unité d'aire, de l'aire $S$ du domaine $D$.
\end{enumerate}
\end{minipage}\hfill
\begin{minipage}{0.43\linewidth}
\psset{unit=0.5cm}
\begin{pspicture*}(-3,-9)(5,3)
\psline{->}(0,0)(1,0)\psline{->}(0,0)(0,1)\uput[d](1,0){\scriptsize $\vect{\imath}$}
\uput[l](0,1){\scriptsize $\vect{\jmath}$}
\psaxes[Dx=10,Dy=20]{->}(0,0)(-3,-9)(5,3)
\psplot[plotpoints=2000,linewidth=1.pt,linecolor=red]{-3}{6}{x dup mul 4 sub 2.71828 x exp div}
\pscustom[fillstyle=hlines,linecolor=red]{\psplot[plotpoints=2000,linewidth=1.2pt,linecolor=red]{-2}{0}{x 2 exp  4 sub 2.71828 x exp div}\psline[linewidth=0.2pt](0,-4)(0,0)(-2,0)}\uput[dl](0,0){\small O}
\uput[r](-2,2){\red $\mathcal{C}_f$}\rput(-1,-3){$D$}
\end{pspicture*}
\end{minipage}

\bigskip

\end{document}
